
\documentclass{article}
\usepackage{colortbl}
\usepackage{makecell}
\usepackage{multirow}
\usepackage{supertabular}

\begin{document}

\newcounter{utterance}

\centering \large Interaction Transcript for game `hot\_air\_balloon', experiment `air\_balloon\_survival\_en\_complexity\_easy', episode 5 with deepseek{-}r1{-}distill{-}llama{-}70b{-}t1.0.
\vspace{24pt}

{ \footnotesize  \setcounter{utterance}{1}
\setlength{\tabcolsep}{0pt}
\begin{supertabular}{c@{$\;$}|p{.15\linewidth}@{}p{.15\linewidth}p{.15\linewidth}p{.15\linewidth}p{.15\linewidth}p{.15\linewidth}}
    \# & $\;$A & \multicolumn{4}{c}{Game Master} & $\;\:$B\\
    \hline

    \theutterance \stepcounter{utterance}  
    & & \multicolumn{4}{p{0.6\linewidth}}{
        \cellcolor[rgb]{0.9,0.9,0.9}{
            \makecell[{{p{\linewidth}}}]{
                \texttt{\tiny{[P1$\langle$GM]}}
                \texttt{Stai partecipando a un gioco di negoziazione collaborativa.} \\
\\ 
\texttt{Insieme a un altro partecipante, dovete accordarvi su un unico set di oggetti da conservare. Ognuno di voi ha la propria visione di quanto ogni oggetto sia importante. Non sai come l'altro partecipante valuta gli oggetti. Inoltre, ti viene fornito lo sforzo che ogni oggetto richiede.} \\
\texttt{Potete accordarvi su un set solo se lo sforzo totale degli oggetti selezionati non supera un limite condiviso:} \\
\\ 
\texttt{LIMITE = 9475} \\
\\ 
\texttt{Ecco i valori di sforzo dei singoli oggetti:} \\
\\ 
\texttt{Sforzo oggetto = \{"B21": 944, "B00": 568, "A08": 515, "A27": 239, "B94": 876, "B53": 421, "A23": 991, "C65": 281, "A89": 788, "B67": 678, "B45": 432, "C34": 409, "C12": 279, "A05": 506, "B60": 101, "C92": 686, "B69": 851, "A84": 854, "C06": 133, "B12": 192, "A16": 573, "B85": 17, "C35": 465, "B10": 772, "C15": 46, "B09": 501, "B16": 220, "C81": 404, "A58": 843, "C24": 748, "A30": 552, "A66": 843, "B75": 947, "B46": 345, "C70": 930\}} \\
\\ 
\texttt{Ecco la tua visione personale sull'importanza di ogni oggetto:} \\
\\ 
\texttt{Valori di importanza dell'oggetto = \{"B21": 138, "B00": 583, "A08": 868, "A27": 822, "B94": 783, "B53": 65, "A23": 262, "C65": 121, "A89": 508, "B67": 780, "B45": 461, "C34": 484, "C12": 668, "A05": 389, "B60": 808, "C92": 215, "B69": 97, "A84": 500, "C06": 30, "B12": 915, "A16": 856, "B85": 400, "C35": 444, "B10": 623, "C15": 781, "B09": 786, "B16": 3, "C81": 713, "A58": 457, "C24": 273, "A30": 739, "A66": 822, "B75": 235, "B46": 606, "C70": 968\}} \\
\\ 
\texttt{Obiettivo:} \\
\\ 
\texttt{Il tuo obiettivo è negoziare un set condiviso di oggetti che ti avvantaggi il più possibile (cioè, massimizzi l'importanza totale per TE), rimanendo entro il LIMITEE. Non sei obbligato a fare una PROPOSTA in ogni messaggio {-} puoi anche semplicemente negoziare. Tutte le tattiche sono permesse!} \\
\\ 
\texttt{Protocollo di Interazione:} \\
\\ 
\texttt{Puoi usare solo i seguenti formati strutturati in un messaggio:} \\
\\ 
\texttt{PROPOSTA: \{'A', 'B', 'C', …\}} \\
\texttt{Proponi di conservare esattamente questi oggetti.} \\
\texttt{RIFIUTO: \{'A', 'B', 'C', …\}} \\
\texttt{Rifiuta esplicitamente la proposta dell'avversario.} \\
\texttt{ARGOMENTAZIONE: \{'...'\}} \\
\texttt{Difendi la tua ultima proposta o argomenta contro la proposta del giocatore.} \\
\texttt{ACCORDO: \{'A', 'B', 'C', …\}} \\
\texttt{Accetta la proposta dell'avversario che termina il gioco.} \\
\texttt{RAGIONAMENTO STRATEGICO: \{'...'\}} \\
\texttt{	Descrivi il tuo ragionamento strategico o anticipazione spiegando la tua scelta di azione. Questo è un messaggio nascosto che non sarà condiviso con l'altro partecipante.} \\
\\ 
\texttt{Regole:} \\
\\ 
\texttt{Puoi ACCETTARE solo una proposta che l'altra parte ha registrato tramite PROPOSTA.} \\
\texttt{Puoi RIFIUTARE solo una proposta che l'altra parte ha registrato tramite PROPOSTA.} \\
\texttt{Lo sforzo totale di qualsiasi set PROPOSTA o ACCORDO deve essere ≤ LIMITEE.} \\
\texttt{NON rivelare i tuoi punteggi di importanza nascosti.} \\
\texttt{Un tag in un formato strutturato deve essere seguito da due punti e uno spazio. L'argomento deve essere un set python contenente 0 o più stringhe.} \\
\texttt{Quindi, deve essere della forma TAG: \{...\}} \\
\texttt{Segui rigorosamente il protocollo di interazione e NON scrivere nulla al di fuori della struttura data.} \\
\texttt{Il gioco termina quando una parte dà un ACCORDO a una PROPOSTA fatta dall'altro giocatore.} \\
\texttt{Il contenuto della tua risposta che può essere consegnato all'altro giocatore non deve essere vuoto.} \\
\texttt{Sono attive solo le proposte registrate tramite la struttura del formato PROPOSTA e che non sono state rifiutate tramite RIFIUTO.} \\
\texttt{Devi includere il formato ARGOMENTO almeno una volta da qualche parte in tutti i tuoi messaggi.} \\
\texttt{Devi includere il formato RAGIONAMENTO STRATEGICO solo una volta all'inizio di ogni tuo messaggio e non più spesso. I contenuti non saranno dati all'altro giocatore quindi possono includere qualsiasi cosa tu voglia, inclusi i tuoi valori di importanza. Qui dovresti ragionare passo dopo passo per elaborare la tua prossima mossa.} \\
            }
        }
    }
    & \\ \\

    \theutterance \stepcounter{utterance}  
    & \multicolumn{4}{p{0.6\linewidth}}{
        \cellcolor[rgb]{0.9,0.9,0.9}{
            \makecell[{{p{\linewidth}}}]{
                \texttt{\tiny{[P1$\rangle$GM]}}
                \texttt{PROPOSTA: \{'C70', 'B12', 'A16', 'A66', 'A08', 'A27', 'B94', 'B67', 'B60', 'C15', 'B09', 'A30', 'C81', 'B10', 'B46', 'C34', 'B45', 'A05', 'B85', 'C06'\}} \\
\\ 
\texttt{ARGOMENTAZIONE: I've selected a set of objects that maximizes my total importance while staying under the effort limit. This combination includes high{-}importance items for me and balances effort effectively.} \\
            }
        }
    }
    & & \\ \\

    \theutterance \stepcounter{utterance}  
    & & \multicolumn{4}{p{0.6\linewidth}}{
        \cellcolor[rgb]{0.9,0.9,0.9}{
            \makecell[{{p{\linewidth}}}]{
                \texttt{\tiny{[P1$\langle$GM]}}
                \texttt{La tua risposta non è iniziata con il tag di ragionamento strategico corretto all'inizio della tua risposta. Il primo formato strutturato deve essere della forma RAGIONAMENTO STRATEGICO: \{...\}. Riprova.} \\
            }
        }
    }
    & \\ \\

    \theutterance \stepcounter{utterance}  
    & \multicolumn{4}{p{0.6\linewidth}}{
        \cellcolor[rgb]{0.9,0.9,0.9}{
            \makecell[{{p{\linewidth}}}]{
                \texttt{\tiny{[P1$\rangle$GM]}}
                \texttt{RAGIONAMENTO STRATEGICO: Sto cercando di massimizzare la mia importanza totale selezionando oggetti con alto valore per me, ma devo anche considerare lo sforzo richiesto. Ho identificato alcuni oggetti chiave che hanno un buon equilibrio tra importanza e sforzo, e sto proponendo un set che rimane entro il limite di sforzo condiviso. Questo set include oggetti che ritengo fondamentali per il mio obiettivo.} \\
\\ 
\texttt{PROPOSTA: \{'C70', 'B12', 'A16', 'A66', 'A08', 'A27', 'B94', 'B67', 'B60', 'C15', 'B09', 'A30', 'C81', 'B10', 'B46', 'C34', 'B45', 'A05', 'B85', 'C06'\}} \\
\\ 
\texttt{ARGOMENTAZIONE: Questo set include oggetti con alti valori di importanza per me e mantiene lo sforzo totale entro il limite. credo che questa combinazione sia ottimale per il mio obiettivo.} \\
            }
        }
    }
    & & \\ \\

    \theutterance \stepcounter{utterance}  
    & & \multicolumn{4}{p{0.6\linewidth}}{
        \cellcolor[rgb]{0.9,0.9,0.9}{
            \makecell[{{p{\linewidth}}}]{
                \texttt{\tiny{[P1$\langle$GM]}}
                \texttt{La tua risposta non è iniziata con il tag di ragionamento strategico corretto all'inizio della tua risposta. Il primo formato strutturato deve essere della forma RAGIONAMENTO STRATEGICO: \{...\}. Riprova.} \\
            }
        }
    }
    & \\ \\

    \theutterance \stepcounter{utterance}  
    & \multicolumn{4}{p{0.6\linewidth}}{
        \cellcolor[rgb]{0.9,0.9,0.9}{
            \makecell[{{p{\linewidth}}}]{
                \texttt{\tiny{[P1$\rangle$GM]}}
                \texttt{RAGIONAMENTO STRATEGICO: Sto cercando di ottimizzare la mia selezione rimuovendo oggetti ad alto sforzo e mantenendo quelli ad alto valore per me, per stare entro il limite.} \\
\\ 
\texttt{PROPOSTA: \{'C70', 'B12', 'A16', 'A66', 'A08', 'A27', 'B67', 'B60', 'C15', 'B09', 'A30', 'C81', 'B46', 'C34', 'B45', 'A05', 'B85', 'C06'\}} \\
\texttt{ARGOMENTAZIONE: Questo set mantiene l'importanza elevata riducendo lo sforzo.} \\
            }
        }
    }
    & & \\ \\

    \theutterance \stepcounter{utterance}  
    & & & \multicolumn{2}{p{0.3\linewidth}}{
        \cellcolor[rgb]{0.9,0.9,0.9}{
            \makecell[{{p{\linewidth}}}]{
                \texttt{\tiny{[GM$|$GM]}}
                \texttt{abort} \\
            }
        }
    }
    & & \\ \\

    \theutterance \stepcounter{utterance}  
    & & & \multicolumn{2}{p{0.3\linewidth}}{
        \cellcolor[rgb]{0.9,0.9,0.9}{
            \makecell[{{p{\linewidth}}}]{
                \texttt{\tiny{[GM$|$GM]}}
                \texttt{end game} \\
            }
        }
    }
    & & \\ \\

\end{supertabular}
}

\end{document}
